\documentclass[12pt]{article}

%%%%%%%%%%%%%%%%%%%%%%% Don't change anything in here. This space is called the preamble, it is where you tell the computer to load the proper LaTeX packages to perform the math and formatting desired. 

\usepackage{physics} 
\usepackage{siunitx} 
\usepackage{enumerate} 
\usepackage{pgfplots}
\usepackage{pgfplotstable}
\usepackage{tikz,pgfplots}
\usepackage{amsmath}  %I added this so that you can use the align tool for equations!
\usepackage{wasysym} %This package allows you to put emojis in your paper!!!!
	%wasysym: \smiley{} \frownie{} see http://milde.users.sourceforge.net/LUCR/Math/mathpackages/wasysym-symbols.pdf for list of most symbols available in this package
	
\usepackage{geometry}
 \geometry{
 a4paper,
 total={170mm,257mm},
 left=20mm,
 top=20mm,
 }

\pgfplotsset{compat=1.14}
%%%%%%%%%%%%%%%%%%%%%%%%% Again, Don't change anything Above %%%%%%%%%%%%%%%%%%%%

\begin{document}

\title{Lab Report Template for CHS Physics 211 v2.0} %Title should be concise and to the point  
\author{Put your name and your partner's names here} %Your name first


\date{\today}  % This will automatically put today's date in the report

\maketitle  %this command makes the title

%%%% to use this template, please copy-paste the entire thing into a new document and save it so you have it!

%%%%% If you want to omit something in this lab, place a % sign to the left of it and it won't show up on the lab, like this line!

\begin{abstract} 
	The abstract should be a short 3-5 sentence paragraph. In it, you should state the hypothesis, what you did, and what you found. The abstract is meant to be a very short summary of your paper to follow. It is a good suggestion to write the abstract last in your report after you have written everything else. This will allow you to best summarize your work. 
\end{abstract}

\section{Introduction}
	 
    This template is designed to be a guide for completing lab reports for Physics 211. Please note that although I have provided this template, you are not under any obligation to use it. It simply provides a nice general design and tips for writing an exceptional physics laboratory report. As you use this template, please make a copy of this code and copy/paste it into a new document so that you will always have access to this original version to look at. It would lso be a good idea to have a printed copy of the pdf template and code to reference when needed. In addition, the formatting for various things like data tables, graphs, citations, etc. is in here for convenience. Please read through the code provided carefully as there are lots of things in \% that you can't see on the generated pdf (like how to include pictures, equations, and LOTS of notes on how to do everything). What follows is a general description of what each section of the lab report should contain. Finally, please note that there are equations and useful latex code at the bottom of this report.\\
    
    The introduction is usually a  paragraph or two. You should start off by stating the problem of the experiment and why the problem is worth studying. If there is any historical or background knowledge required to understand your report, you should state it here. Finally, you should state your hypothesis as an if-then statement. \cite{1}
    
    %note the cite above. That is an in text citation. By putting citations in your bibliography, you can cite in text simply using the \cite{} command and putting the number of the citation in the {} 
    

    
\section{Experimental Setup}
	In this section, you should explain your experiment in detail. The usual method is that you first explain what you did, then you explain how you did it. I.e. first explain the data you are trying to collect. Then explain (citing relevant equations as necessary) why this data is necessary in order to prove/disprove your hypothesis. Finally, explain how you collected the data, citing both the materials used as well as stating their purpose and what they are used to collect or to do. You can break up this section into two sections if you wish. The first would be materials and the second would be procedure. 

%%%%%%%%%%Here is a sample bulleted list that is often useful for the materials section. You can also number or letter your materials to easily refer to them later on
\begin{description}
  \item[$\bullet$] 61cm ramp
  \item[$\bullet$] stack of textbooks
  \item[$\bullet$] ball
  \item[$\bullet$] meter stick
  \item[$\bullet$] tape
  \item[$\bullet$] timer
  \item[$\bullet$] Logger Pro
  \item[$\bullet$] motion sensor
\end{description}
%%%%%%%%%%%%%%%%%%%%%%%%%%%%%%%%%%%%%%%%%%%%%%


\section{Data}

In this section you should explain your results. You should not interpret the results in this section, that is for the analysis/conclusion portion of the lab. Again, you are merely stating the results of your experiment here. Please avoid gigantic tables with every data point in existence that no one will ever read. Use data tables wisely to summarize large amounts of data or to highlight important trends/results. If you want to display large amounts of data, often a graph/histogram, pie chart, etc. is the wiser choice. 

%%%%%Here is a nice sample data table%%%%%%%%%%
\begin{center}
\begin{tabular}{|c|c|c|c|c|} %this specifies how many columns to make. For example. if you only wanted three columns, it would look like this : {|c|c|c|}
\hline
Angle of incline & Trial 1 (Sec.) & Trial 2 (Sec.) & Trial 3 (Sec.) & Average Time (Sec.)\\[0.5ex] %Title for each column
\hline\hline
$2^{\circ}$&$2.615$&$2.73$5&$2.585$&$2.645$\\ %data for row 1
\hline
$4^{\circ}$&$1.8$&$1.875$&$1.685$&$1.787$\\ %data for row 2
\hline
$6^{\circ}$&$1.445$&$1.39$&$1.44$&$1.425$\\ %data for row 3
\hline
8$^{\circ}$&1.195&1.31&1.1&1.202\\ %data for row 4
\hline
\end{tabular}
\end{center} 
 
% When making a table, make sure every data point is separated by $ to signify a different cell in your table. To end a row, use the \\ to move to the next row. 
%%%%%%%%%%%%%%%%%%%%%%%%%%%%%%%%%%%%%%%%%%%%%%%%%%


%%%%%%%%%%    Using Graphs/pictures in LaTeX    %%%%%%%%%%%%%%%%%%%%%%%%%%
%In this section you will find the code for inserting a graph into LaTeX. My suggestion is to use excel to create your graph, so that you can include trend lines, error bars (set to 2 sigma standard deviation or to a set amount, e.g. 0.2s to include roughly for human error effects due to timing. You can also use excel to plot experimental data (with error bars) and theoretical data, which allows you to easily see if the theoretical data falls within the region of acceptable error. 

%Once your graph is created, take a screen shot using print screen (prt scn on most computers). Paste into paint and crop out your graph and save as a jpg or png file. Title your picture as something simple to reference later, for example "graph1"

%You will then click the "project" tab in the upper left-hand corner of overleaf, click the drop-down arrow "main.tex" and click upload to upload your picture. From there, use the code below to insert the picture into your report. 

%technique #1: Basic way
%\includegraphics[width=\textwidth]{YOUR-FILE-NAME-HERE}

%technique #2: slightly more advanced way which allows formatting, sizing, cen

%\begin{figure}[h!]
 % \caption{A picture of a gull.}
  %\centering
  %\includegraphics[width=0.5\textwidth]{Image File Name}
%\end{figure}

%%%%%%%%%%%%%%%%%%%%%%%%%%%%%%%%%%%%%%%%%%%%%%%%%%%%%%%%%%%%%%%%%%%%%%%%%%%%%%%%





%%%%%%%%%%%%%%%%%%%%%%%%%%%%%%%%%%%%%%%%%%%%%%%%%%%%%%%%%%%%%%%%%%%%%%%%%%%%%%%%%
%Here are some sample equations showing how to use the align tool to number and align math in your report

\begin{align} %Note that in align, you are in "math mode" and thus you don't need the $ between your math
K &= U_g \\
mgh &=\frac{1}{2}mv^2 \\
(10)(.225) &=\frac{1}{2}v^2 
onumber \\  %by using 
onumber, you can take off the label on the side of an equation, which is nice so that steps of math with work aren't numbered. It is entirely your choice if you want to use this or not. 
4.5 &=v^2 
onumber \\
2.12 \frac{m}{s} &=v 
\end{align}
%The align tool allows you to number your equations, which is super useful for referencing math in your work. I.e. you could say "see equation (1) for details." 

%Note: the \\ is to make a new line. Whatever you want aligned, just use & to the left, for example see the = above. Finally, don't use \\ on the last line or it will create a blank bottom equation
%%%%%%%%%%%%%%%%%%%%%%%%%%%%%%%%%%%%%%%%%%%%%%%%%%%%%%%%%%%%%%%%%%%%%%%%%%%%%%%%%%%



%~~~~may the f=ma be with you!~~~~~~~~Mr. C 



%%%%%%%%%%% Here is how to include a picture in your report   %%%%%%%%%%%%%%%%%%%%%
% \includegraphics[width=.3\textwidth]{Picture13}
% In this example, Picture2.png is the name of the picture file. For every picture that you want, you first have to save the picture to your computer, then give it a name like Picture2.png 
%%%%%%%%%%%%%%%%%%%%%%%%%%%%%%%%%%%%%%%%%%%%%%%%%%%%%%%%%%%%%%%%%%%%%%%%%%%%%%%%%%%%



\section{Analysis}
	This is the most important part of the lab report; it is where you
analyze the data. In this section you will interpret your results. You need to look at your data and decide if the hypothesis was supported or contradicted by your data.  \\ %Note: the \\ gives you a line of space
    
    Your discussion should include the following at a minimum. [1] What is the relationship between your measurements and your final results? [2] What trends were observable? [3] What can you conclude from the graphs that you made? [4] How did the independent variables affect the dependent variables? (For example, did an increase in a given measured (independent) variable result in an increase or decrease in the associated calculated (dependent) variable?) \\ 

Then describe how your experimental results substantiate/agree with the theory. (This is not a single statement that your results agree or disagree with theory.) When comparison values are available, discuss the agreement using either uncertainty and/or percent differences. This leads into the discussion of the sources of error. Your discussion should include the calculation of averages and standard deviation to be able to describe precision of experiment. All data points should be plotted $\pm$ two standard deviations and compared with theoretical data to interpret accuracy. It is ok to say that your results were inconclusive. It is important to cite all possible sources of error and state specifically how you believed they affected the collected data. If you get a result or an uncertainty that is ridiculous (or just really big/small), show that you have noticed and thought about it, not just copied a number from your calculator and moved on. \\

For example, when rolling a ball down a ramp, you may not have taken in account the effects of rolling friction or the fact that some gravitational potential energy is converted into rotational kinetic energy. Both of these would cause the overall time for the ball to roll down the ramp to increase.


\section{Conclusion}
    The conclusion should connect to the introduction and re-state the relevance and importance of the experiment. Its a nice touch to sometimes make historical connections in this part of the report as well. It is always good to end on a note stating the importance of your findings, the connections to other topics in physics and science, and opportunities for future extensions/research/experiments in the subject. Remember to report your results with correct units and uncertainties, for example $g=9.7 \pm 0.2 m \cdotp s^{-2}$. 
    
%%%%%%%%%%%%%%%%%%%%%%%%%%%%%%%%%%%%%%%%%%%%%%%%%%%%%%%%%%%%%%%%%%%%%%%%%%%%%%%%%%%%%
%%%% This is the Bibliography where you will cite your sources used in the paper %%%%

\begin{thebibliography}{0}

	%Each item starts with a \bibitem{} command and the details thereafter.
	
	\bibitem{1} Cite your first source here
	\bibitem{2} Cite another source
	
    %%% The 1,2 etc. are used to cite in text. See up in the intro for an example
    %%% When you want to cite in your cite, type in \cite{} wherever you want
\end{thebibliography}
    
%%%%%%%%%%%%%%%%%%%%%%%%%%%%%%%%%%%%%%%%%%%%%%%%%%%%%%%%%%%%%%%%%%%%%%%%%%%%%%%%%%%%    


%%%%%%%%%%%%%%%%%%%%%%%%%%%%%%%%%%%%%%%%%%%%%%%%%%%%%%%%%%%%%%%%%%%%%%%%%%%%%%%%%%%%%

% Here are some equations and other useful things, feel free to copy and paste into your lab report :) %

% Note: It is super easy to look up equations/constants already formatted in LaTeX online. Here are a few websites I like to use:

% LaTeX Tutorial: http://pages.physics.cornell.edu/sps/pages/resources/latex.html
	%NOTE: This contains both pdf and text (code) of each document, which includes guides, lab 			report templates, and lots of other good stuff!

% LaTeX Cheat Sheet: http://wch.github.io/latexsheet/

% Equations: http://www.equationsheet.com/sheets/Equations-5.html

% Constants, symbols, letters, etc:  http://www.rpi.edu/dept/arc/training/latex/LaTeX_symbols.pdf

% AP Physics Calculus Reference Table: https://secure-media.collegeboard.org/digitalServices/pdf/ap/physics-c-tables-and-equations-list.pdf

% AP Physics Trig Reference Table: https://secure-media.collegeboard.org/digitalServices/pdf/ap/ap-physics-1-equations-table.pdf

% Regents Physics Reference Table: http://www.p12.nysed.gov/assessment/reftable/physics-rt/physics06tbl.pdf

% Here are three good sources to use other then this for considering how to write a good lab report. Much of this guide was gleaned from them
	
    %http://web.mit.edu/8.13/www/Samplepaper/simple-zipped/simple-paper.pdf 				-MIT Physics Lab (really, really good template!) 
    
    % http://pages.physics.cornell.edu/sps/pages/resources/LatexSession/Exercises/LabReport.pdf -Cornell Template
    % https://engineering.purdue.edu/ME588/LabManual/report_format.pdf 							-Purdue Engineering Template 
    % http://physics.columbia.edu/files/physics/content/1291_report_format_and_example.pdf 		-Columbia University Template
    % https://www.baylor.edu/physics/doc.php/110769.pdf											-Baylor University Template
    % www.nd.edu/~hgberry/Fall2012/Guidelines.docx												-Notre Dame Template
    %http://www.esf.edu/iq/colloquium/documents/LabReportnotes.pdf								-SUNY ESF Template
    %http://writing.engr.psu.edu/workbooks/laboratory.html										-Virginia Tech Template
    %https://projects.ncsu.edu/labwrite/index_labwrite.htm										-SUPER in-depth guide to writing lab reports
    
    %https://gist.github.com/dcernst/1827406													-Template for completing Math homework in LaTeX
    %https://joshldavis.com/2014/02/12/doing-your-homework-in-latex/							-More about math homework in LaTeX
    
    
    
%Purdue University Online Writing Lab-OWL: https://owl.english.purdue.edu/ Use this to generate citations!

%Basically, if you get stuck, just google "Latex ______" for whatever you need and look through the LaTeX stackexchange or wiki article to find and copy/paste what you need





%	Here are some common equations we use in class. I will continue to update as we continue throughout the year. I will attempt to organize by the order we learn the topics from oldest at the top to newest at the bottom. Go ahead and copy/paste as needed in your report

%%%%%%%%%%%%%%%%%%%%%%%%%% Basic Calculus %%%%%%%%%%%%%%%%%%%%%%%%%%%%%%%%%%%%%

%		   $$\frac{\mathrm{d}}{\mathrm{d}x}C=0$$
%           $$\frac{\mathrm{d}}{\mathrm{d}x}Cx=C$$
%           $$\frac{\mathrm{d}}{\mathrm{d}x}x=1$$           
%           $$\frac{\mathrm{d}}{\mathrm{d}x}x^n = nx^{n-1} $$  	-power rule
% 		   $$\frac{\mathrm{d}}{\mathrm{d}x}fg= fg'+f'g$$		-product rule
%           $$\frac{\mathrm{d}}{\mathrm{d}x}f(g(x))=f'(g(x))g'(x)$$	-chain rule
%           $$\frac{\mathrm{d}}{\mathrm{d}x} \sin{x} = \cos{x}$$	
%           $$\frac{\mathrm{d}}{\mathrm{d}x} \cos{x} = -\sin{x}$$
           
%           $$\int k \mathrm{d}x = kx+C$$		-integral of a constant
%           $$\int x^n \mathrm{d}x= \frac{1}{n+1}x^{n+1}+C$$	-power rule for integrals
%           $$\int \cos{u}\mathrm{d}u = \sin{u} + C$$	
%           $$\int \sin{u}\mathrm{d}u = -cos{u} + C$$	

% see https://reu.dimacs.rutgers.edu/Symbols.pdf for a nice list of math LaTeX symbols

%%%%%%%%%%%%%%%%%%%%%%%%%%%%  Kinematics %%%%%%%%%%%%%%%%%%%%%%%%%%%%%%%%%%%%%%%

%			$$\bar{v}=\frac{d}{t}$$ 								-average speed
%			$$v=\frac{\mathrm{d}x}{\mathrm{d}t}$$					-instantaneous velocity definition
%			$$a=\frac{\mathrm{d}v}{\mathrm{d}t}$$					-instantaneous acceleration definition
%			if acceleration is constant, then:
%				$${x_f}={x_i}+{v_i}t+\frac{1}{2}a{t^2}$$  			-free-fall equation
%				$${v_f}={v_i}+at$$									-find new velocity
%				$${{v_f}^2}={{v_i}^2}+2a({x_f}-{x_i})$$				-equation without time

%%%%%%%%%%%%%%%%%%%%%%%%%%%%% Newton's Laws %%%%%%%%%%%%%%%%%%%%%%%%%%%%%%%%%%%%%%

%			$$F_{net}=ma=m\frac{\mathrm{d}v}{\mathrm{d}t}=m\frac{\mathrm{d}^{2}x}{\mathrm{d}{t^2}}$$	-Newt's 2nd Law
%			$$F_f=\mu F_n$$											-Friction 
%			$$w=mg$$												-Weight

%%%%%%%%%%%%%%%%%%%%%%%%%%%%%%%% Work, Power, Energy %%%%%%%%%%%%%%%%%%%%%%%%%%%%%%%%%%%

%			$$W=\Delta E = \int F dx = Fd  \hspace{3pt} \text{(if F constant)} $$ 	-Work-Energy Theorem
%			$$Power=\frac{\mathrm{d}E}{\mathrm{d}t}=\frac{\mathrm{d}W}{\mathrm{d}t} = \frac{Fd}{t} = F \bar{v}$$		-Power
%			$$K=\frac{1}{2}mv^2$$													-Kinetic Energy
%			$$U_g=mgh$$																-Gravitational Potential Energy
%			$$U_e=\frac{1}{2}kx^2$$													-Spring Potential Energy
%			$$F_e=kx$$																-Hooke's Law
%			$$F=-\frac{\mathrm{d}U}{\mathrm{d}x}$$									-Force is derivative of Potential

%%%%%%%%%%%%%%%%%%%%%%%%%%%%%%%%%%%%% Momentum, Center of Mass %%%%%%%%%%%%%%%%%%%%%%%%%%%%%%%%%%%

% 			$$p=mv$$								-definition of momentum
%			$$F=\frac{\mathrm{d}p}{\mathrm{d}t}
%			$$ft=\Delta{p}$$ 						-trig version of impulse momentum theorem
%			$$J=\int F \mathrm{d}t = \Delta{p}$$	-calc version of impulse momentum theorem
%			$$p_{before}=p_{after}$$				-Conservation of Momentum
%			$$X_{c.o.m}=\frac{\Sigma x_i m_i}{M}$$  -x-coordinate of Center of Mass
%       	$$Y_{c.o.m}=\frac{\Sigma y_i m_i}{M}$$  -y-coordinate of Center of Mass


%%%%%%%%%%%%%%%%%%%%%%%%%%%%%%%%%%% Rotational Kinematics %%%%%%%%%%%%%%%%%%%%%%%%%%%%%%%%%%%%%%%

%		$$\omega=\frac{\mathrm{d}\theta}{\mathrm{d}t}$$ 			-definition of angular speed (rad/sec)
%		$$\alpha=\frac{\mathrm{d}\omega}{\mathrm{d}t}=\frac{\mathrm{d}^2\theta}{\mathrm{d}t^2}$$ 			-definition of angular acceleration
%		$$v=r\omega$$												-angular/linear velocity connection
%		$$S=r \theta$$												-arc length/angle connection
%		if angular acceleration is constant, then:
%			$$\omega_f=\omega_i+\alpha t$$   	-find angular velocity
%			$$\theta=\theta_i + \omega_i t + \frac{1}{2} \alpha t^2$$ 		-angular "free-fall" equation
%			$${\omega_f}^2 = {\omega_i}^2 + 2\alpha(\theta_f - \theta_i)$$ 	-angular equation w/o time

%%%%%%%%%%%%%%%%%%%%%%%%%%%%%%%% Rotational Dynamics + Gravitation %%%%%%%%%%%%%%%%%%%%%%%%%%%%%%%%%%%%%%%%%%%

%		$$\tau=r \times F = rF\sin{\theta}$$			-definition of torque
%		$$\tau = I \alpha$$								-Newt's 2nd Law for Rotation
%		$$a_c = v^{2}/r = {\omega^2}r$$					-centripetal acceleration
%		$$F_c=ma_c = m{\omega^2}r$$						-centripetal force
%		$$I=\int r^2 \mathrm{d}m = \Sigma m r^2$$		-Calculate Moment of Inertia of an Object
%		$$K_r=\frac{1}{2}I{\omega^2}$$					-Rotational Kinetic Energy
%		$$L= I\omega = r \times p = rp\sin{\theta}$$	-Angular Momentum
%		$$L_{before}=L_{after}$$						-Conservation of Angular Momentum

%		$$F_g = \frac{\text{G}m_1 m_2}{r^2}$$			-Newton's Law of Universal Gravitation
%		$$U_g= -\frac{\text{G}m_1 m_2}{r}$$				-Gravitational Potential Energy

%%%%%%%%%%%%%%%%%%%%%%%%% Vibrations, Simple Harmonic Motion, Sound %%%%%%%%%%%%%%%%%%%%%%%%%%%%%%	

%		$$v=f\lambda$$								-speed of a wave
%		$4T=\frac{2\pi}{\omega}=\frac{1}{f}$$		-period of a wave
%		$$x(t)=x_{max}\cos{(\omega t + \phi)}$$	    -wave equation
%		$$T_s=2\pi\sqrt{\frac{m}{k}}$$				-Period of an oscillating spring
%		$$T_p=2\pi\sqrt{\frac{l}{g}}$$				-Period of a Pendulum
%		$$\frac{\mathrm{d}^2 \smiley{}}{{dt}^2} - {\omega}^2 \smiley{} = 0$$ 	-General Simple Harmonic Motion Equation
%		
		
%%%%%%%%%%%%%%%%%%%%%%%%%%%%%%%%%%%%%% Fluid Mechanics %%%%%%%%%%%%%%%%%%%%%%%%%%%%%%%%%%%%%%%%%%%

%		$$\rho=m/V$$						-Density
%		$$P=F/A$$							-Definition of Pressure
%		$$P=P_i+\rho gh	$$					-Pressure change as a function of depth
%		$$\rho_1A_1v_1=\rho_2A_2v_2$$		-Continuity Equation for fluid flow
%		$$F_b=\rho gV_{displaced}$$			-Archimedes Principle
%		$$P_1+\frac{1}{2}\rho v_1^2 +\rho gh_1 = P_2+\frac{1}{2}\rho v_2^2 +\rho gh_2 $$ -Bernoulli Equation
%		

%%%%%%%%%%%%%%%%%%%%%%%%%%%%%%%%%%%%%% Thermodynamics %%%%%%%%%%%%%%%%%%%%%%%%%%%%%%%%%%%%%%%%%%%%%

%		$$\frac{\Delta Q}{\Delta t} = \frac{k A \Delta T}{l}$$ 		-Conduction Equation
%		$$Q = mc \Delta T$$											-Heat Flow Equation (sensible heat)
%		$$Q=mL_f$$													-Latent Heat of Fusion
%		$$Q=mL_v$$													-Latent Heat of Vaporization
%		$$ \Delta l = \alpha l_0 \Delta T$$							-Length expansion
%		$$ \Delta A = 2 \alpha A_0 \Delta T$$						-Area expansion
%		$$ \Delta V = 3 \alpha V_0 \Delta T$$						-Volume Expansion
%       $$PV=Nk_BT$$												-Ideal Gas Law
%		$$K=\frac{n}{2}k_BT$$										-Equipartition Theorem
			%if Ideal Gas:
            	% then $$K=\frac{3}{2}k_BT$$
%		$$W=-\int P \mathrm{d}V										-Thermodynamic Work on a gas (calc)
%		$$W=-P \DeltaV$$											-Thermodynamic Work on a gas Equation-trig
%		$$\Delta U= Q+W$$											-1st Law of Thermodynamics
%		$$\epsilon_{real} = 1- \frac{Q_c}{Q_h}$$					-Real Efficiency 
%       $$\epsilon_{theory} = 1- \frac{T_c}{T_h}$$					-Theoretical "Carnot" efficiency


%%%%%%%%%%%%%%%%%%%%%%%%%%%%%%%%%%%%%%% Misc. Useful Things %%%%%%%%%%%%%%%%%%%%%%%%%%%%%%%%%%%%%%%
% 
% \framebox{box}  This puts a box around something. Good for showing something important
% to quote someone, use the following template:
	%\begin{quotation}
	%``You miss 100\% of the shots you never take'' %note that quotes are formatted this way in LaTeX
	%-Michael Scott
	%\end{quotation}



% https://physics.info/equations/ is a good source for most common equations found in introductory physics (trig and calc based)




\end{document}